\hspace{3mm}

\centerline{\bf \Huge Алгоритмы без обратной связи}

\begin{flushright}
    \scriptsize{Честно я отсюда ни одну задачу не решал)}
\end{flushright}


\problem{1} В одном из 1000 окопов, расположенных в ряд, спрятался робот. Пушка может од- ним выстрелом накрыть любой окоп. В каждом промежутке между выстрелами ро- бот (если уцелел) обязательно перебегает в соседний окоп (быть может, только что обстрелянный). Сможет ли пушка наверняка накрыть робота?

\problem{2} Левша и невидимая блоха на плоскости играют, ходя по очереди. Очередным ходом Левша проводит прямую, а блоха совершает прыжок длины 1, не пересекающий ни одной прямой. Если таких прыжков нет, блоха проигрывает. Может ли Левша выиграть, как бы не играла блоха?

\problem{3} Ёжик стоит в левой нижней клетке поля 8 × 8. А в какой-то другой клетке пасётся лошадка. На поле стоит туман, ничего не видно, но Ёжику надо найти Лошадку. Лошадка каждую минуту переходит на соседнюю по стороне клетку и говорит, куда она перешла (влево, вправо, вверх или вниз). Ёжик тоже может сделать шаг в одну из соседних по стороне или диагонали клеток, как только услышит Лошадку. Ёжик найдёт Лошадку, если окажется с ней на одной клетке. Что же делать Ёжику?

\problem{4} Назовём лабиринтом доску 8 × 8, где между некоторыми полями вставлены перего- родки. По команде ВПРАВО ладья смещается на одно поле вправо или, если справа край доски или перегородка, остается на месте; аналогично выполняются команды ВЛЕВО, ВВЕРХ и ВНИЗ. Петя пишет программу — конечную последовательность указанных команд, и дает её Васе, после чего Вася выбирает лабиринт и помещает в него ладью на любое поле. Докажите, что Петя может написать такую программу, что ладья обойдет все доступные поля в лабиринте при любом выборе Васи.

\problem{5} Капитан Врунгель в своей каюте разложил перетасованную колоду из 52 карт по кругу, оставив одно место свободным. Матрос Фукс с палубы, не отходя от штурва- ла и не зная начальной раскладки, называет карту. Если эта карта лежит рядом со свободным местом, Врунгель её туда передвигает, не сообщая Фуксу. Иначе ничего не происходит. Потом Фукс называет ещё одну карту, и так сколько угодно раз, пока он не скажет «стоп».

(а) Может ли Фукс добиться того,чтобы после «стопа» каждая карта наверняка оказалась не там, где была вначале?

(б) Можетли Фукс добиться того,чтобы после «стопа» рядом со свободным местом наверняка не было туза пик?

\problem{6.} На бесконечной в обе стороны улице стоит отделение полиции, из которого сбежал подозреваемый. Максимальная скорость полицейского $- 1$, подозреваемого $- v$. Время побега и местоположение подозреваемого полицейскому не известны. Верно ли, что он сможет поймать подозреваемого (оказаться с ним в одной точке), если ему известно, что 
(а) $v = 0,9;$ (б) $v < 1;$ (в) То же условие, $v < 1,$ но теперь вместо бесконечной прямой точка с $v$ выходящими из неё бесконечными лучами.